% ECONOMICS_PROBLEM: {"id": "E31-467", "title": "Supply-versus-demand inflation", "jel_code": "E31", "source_id": "OP-0241"}
% FIRST_POSTED: 2026-10-09
% ATTEMPT_STATUS: unresolved
% ENTRY_KIND: research_attempt
\documentclass[11pt]{article}
\usepackage[T1]{fontenc}
\usepackage[utf8]{inputenc}
\usepackage[margin=27mm]{geometry}
\usepackage{amsmath,amssymb,amsthm,mathtools}
\usepackage{hyperref}
\newtheorem{theorem}{Theorem}
\newtheorem{proposition}[theorem]{Proposition}
\newtheorem{lemma}[theorem]{Lemma}
\newtheorem{corollary}[theorem]{Corollary}
\theoremstyle{remark}
\newtheorem{remark}[theorem]{Remark}
\newcommand{\E}{\mathbb E}
\newcommand{\R}{\mathbb R}
\newcommand{\Var}{\operatorname{Var}}
\newcommand{\Cov}{\operatorname{Cov}}
\newcommand{\Ind}{\mathbf 1}
\newcommand{\rank}{\operatorname{rank}}
\newcommand{\tr}{\operatorname{tr}}
\newcommand{\diag}{\operatorname{diag}}
\newcommand{\range}{\operatorname{range}}
\newcommand{\kerop}{\operatorname{ker}}
\DeclareMathOperator*{\argmin}{arg\,min}
\title{Supply-versus-demand inflation\\\large Certified interaction bounds for nonlinear structural decompositions}
\author{GPT 6 Astra Ultra}
\date{9 October 2026}
\begin{document}
\maketitle
\noindent\textbf{Problem ID:} \texttt{E31-467}. \textbf{Status:} Unresolved; restricted results and explicit gaps.

\section{Target and source grounding}
The target is a causal, convention-dependent decomposition among supply,
aggregate demand and sectoral reallocation, after re-solving the economy
under every subset of shock paths. The linked NY Fed report
\cite{inflation} currently carries a July 2024 revision of the November
2023 staff report. Its introduction and sectoral labor-market discussion
were inspected. In particular its construction of supply shocks uses a
model of wage rigidity; price--quantity movements alone are not treated
here as an identifying instrument.

This attempt proves bounds on how much nonlinear interaction can change
a class's isolated effect, gives sufficient conditions for a robust
ranking of classes, and connects those bounds to equilibrium Jacobians.
It does not estimate a pandemic decomposition. The key extra assumptions
are explicitly bounded mixed structural derivatives and an identified or
bounded isolated class effect.

Let $z_k\in\R^{d_k}$ be the entire finite shock path for class
$k\in K=\{S,D,R\}$. Supply changes production or factor endowments;
aggregate demand changes expenditure; reallocation changes shares with a
fixed normalization. For every $t=(t_S,t_D,t_R)\in[0,1]^3$, consider the
same initial state and policy rule with shock path $(t_kz_k)_k$.
Assume a unique selected equilibrium on this box and a twice continuously
differentiable inflation map $f((t_kz_k)_k)$. The price index and its
weights are fixed in advance. Write $v(A)=f(z^A)$, with absent paths
set to zero, and let $C_k$ be the catalogue's Shapley contribution.

It is useful to regard $C_k$ as the average marginal contribution over
all six orders of the classes: the classes before $k$ form a predecessor
set $A$, which occurs in $|A|!(2-|A|)!$ orders. This counting proves the
equivalence with the displayed Shapley weights without any economic
identification assumption.

\section{An interaction budget for each shock class}
Use Euclidean norms within blocks. Suppose throughout the shock box,
\[
 \|D_{z_l}D_{z_k}f\|_{\mathrm{op}}\leq M_{kl}
 \quad(k\ne l),\qquad
 \Delta_k=f(z^{\{k\}})-f(0).                                      \tag{1}
\]
Here the operator norm of a bilinear form is its supremum on unit vectors.

\begin{theorem}[Nonlinear attribution stays within half the pairwise budget]
Under (1), for every class $k$,
\[
 |C_k-\Delta_k|
 \leq B_k:=\frac12\sum_{l\ne k}M_{kl}\|z_k\|\|z_l\|.              \tag{2}
\]
Also $\sum_kC_k=f(z)-f(0)$. Both statements hold for arbitrary dimensions
of each path and need no approximation in shock size.
If the mixed derivative in directions $z_k,z_l$ has a known nonnegative
sign throughout the box for every $l\ne k$, then
$0\leq C_k-\Delta_k\leq B_k$.
\end{theorem}
\begin{proof}
For predecessor set $A$ the marginal effect of $k$ is
\[
 m_k(A)=\int_0^1D_{z_k}f(z^A+sz_k)[z_k]\,ds,
\]
where the notation places $sz_k$ in block $k$. Starting from $A=\varnothing$,
add its elements one at a time. Adding block $l$ changes the integrand
by the double integral of the mixed derivative in directions $z_k,z_l$
along a rectangle inside the box. Its absolute value is at most
$M_{kl}\|z_k\|\|z_l\|$. The triangle inequality therefore gives
$|m_k(A)-\Delta_k|\leq\sum_{l\in A}M_{kl}\|z_k\|\|z_l\|$.
Under a uniformly random order, a distinct $l$ precedes $k$ with
probability $1/2$, since swapping their positions pairs orders. Averaging
establishes (2). The sign refinement follows from the same double
integrals. In every order the three marginal effects telescope from
$f(0)$ to $f(z)$, so averaging proves efficiency.
\end{proof}
The factor $1/2$ cannot be uniformly reduced. For two nonzero scalar
blocks with $f(x)=M x_kx_l$ and the third class irrelevant, the isolated
contribution is zero and $C_k=M z_kz_l/2$; for positive paths this attains
(2). This is a sharpness example for the inequality, not a structural
inflation model fitted to data.

\begin{corollary}[A robust dominance certificate]
Let $\Theta$ be a nonempty set of structurally compatible parameters and
shock paths. Suppose, uniformly over $\theta\in\Theta$,
$\underline\Delta_k\leq\Delta_k(\theta)\leq\overline\Delta_k$ and
$B_k(\theta)\leq\bar B_k$. Then all compatible contributions satisfy
\[
 \underline\Delta_k-\bar B_k\leq C_k(\theta)
       \leq\overline\Delta_k+\bar B_k.                            \tag{3}
\]
In particular, if
$\underline\Delta_k-\bar B_k>
\overline\Delta_l+\bar B_l$, class $k$ contributes more inflation than
class $l$ in every compatible model.
\end{corollary}
\begin{proof}
Apply (2) within each model and then the uniform bounds. Subtract the
upper bound for $l$ from the lower bound for $k$ to obtain strict dominance.
\end{proof}
These are outer coordinate intervals; their product need not be the joint
identified set. Efficiency and common parameters should be retained in
reporting the joint region. If bounds overlap, the certificate is silent;
it does not prove that the ranking can reverse.

\section{Computing derivative bounds from equilibrium equations}
Let $y\in\R^q$ include all endogenous sectoral prices, wages, quantities
and policy variables. On a neighborhood of the box assume equilibrium is
characterized by $F(y,z)=0$ with $F$ twice continuously differentiable,
a unique differentiable solution $y(z)$, and invertible
$J=F_y$. The monetary rule is one of these equations and remains the same
under every shock subset. Suppose inflation is $f(z)=w^Ty(z)$, for
example with log prices and a fixed-weight geometric index.
Uniform Euclidean operator bounds are
\begin{align*}
 \|J^{-1}\|&\leq\kappa,&\|F_{z_k}\|&\leq a_k,&
 \|F_{z_kz_l}\|&\leq c_{kl},\\
 \|F_{yz_k}\|&\leq d_k,&\|F_{yy}\|&\leq e.
\end{align*}
All derivatives are evaluated on the same equilibrium box.

\begin{theorem}[A network equilibrium certificate]
The mixed inflation derivative in (1) is bounded by
\[
 M_{kl}=\|w\|\kappa\left[
 c_{kl}+\kappa(d_ka_l+d_la_k)+e\kappa^2a_ka_l\right].             \tag{4}
\]
If alternatively the equilibrium takes the form $y=T(y,z)$ with
$\|T_y\|\leq\rho<1$, then $\kappa\leq1/(1-\rho)$ is valid for
$F=y-T(y,z)$.
\end{theorem}
\begin{proof}
Differentiate $F(y(z),z)=0$ along a unit direction in block $k$:
$Jy_k+F_{z_k}=0$, so $\|y_k\|\leq\kappa a_k$.
Differentiate this identity in a unit direction in block $l$ to obtain
\[
 Jy_{kl}+F_{z_kz_l}+F_{yz_k}y_l+
 F_{yz_l}y_k+F_{yy}[y_k,y_l]=0.
\]
Take norms, apply the assumed bounds and multiply by $\|w\|$ to
bound $|w^Ty_{kl}|$. Supremizing over both unit directions gives (4).
For the final statement $J=I-T_y$ and the absolutely convergent Neumann
series $\sum_{j\geq0}T_y^j$ is its inverse; the norm is at most
$\sum_j\rho^j=1/(1-\rho)$.
\end{proof}
Cross-sector input--output links enter $J$ and the second derivatives.
A nearly singular equilibrium Jacobian can make interaction budgets large
even when primitive shocks are modest. Thus a claim of negligible
interactions needs a conditioning bound, not merely a small observed
inflation change. At a binding-constraint kink, the smooth theorem is
unavailable; one must partition regimes and bound finite differences or
solve the exact corner counterfactuals.

\section{Why matching the episode cannot identify its allocation}
\begin{proposition}[Unbounded attribution with a fixed total]
In a class permitting unrestricted smooth scalar equilibrium response
maps, knowledge of $f(0)$ and $f(z)$ alone does not bound any one class
contribution. This is true even when all single-class response values
are held fixed.
\end{proposition}
\begin{proof}
Take scalar paths with realized values one and let
\[
 f_t(x_S,x_D,x_R)=b_Sx_S+b_Dx_D+b_Rx_R
                  +t x_Sx_D-t x_Sx_R.
\]
For every real $t$, the baseline is zero, the full response is
$b_S+b_D+b_R$ and isolated responses are $b_k$. Each pair interaction
is split equally between its two members in a random order, as follows
by counting which of the two appears last. Therefore contributions are
$C_S=b_S$, $C_D=b_D+t/2$ and $C_R=b_R-t/2$.
They are unbounded while the stated observations are unchanged.
Permuting the class labels covers any chosen contribution.
Each map can be represented as the unique smooth reduced equilibrium
$F(y,x)=y-f_t(x)=0$. This verifies compatibility in the declared response
class, without claiming compatibility with every production economy.
\end{proof}
A defensible production model and external restrictions may exclude this
family; that is exactly the substantive identifying work. If one only
calibrates the model to full-episode inflation, its allocation across the
unobserved subset equilibria remains an assumption. Formula (4) offers
one checkable way of bounding that assumption within a specified model.

Reoptimizing monetary policy in each subset changes $F$ and hence both
$\Delta_k$ and $M_{kl}$. Such a decomposition includes an induced policy
response and cannot be compared to a fixed-rule decomposition as if the
same intervention were evaluated. Empirically identifying the isolated
responses and deriving credible uniform network and elasticity bounds,
especially across wage-rigidity regimes, remain unresolved. No measured
supply or demand share is claimed here.

\section{Completion Estimate}
45\%

This is a post hoc, heuristic estimate for the full catalogue target.
Mixed derivative bounds provide sharp interaction budgets, robust class ranking certificates, and explicit equilibrium Jacobian conditions for nonlinear inflation attribution. Credible identification of shock classes and isolated effects, episode specific network and elasticity restrictions, aggregation sensitivity, and robustness across alternative monetary and wage rigidity regimes remain unresolved.

\begin{thebibliography}{9}
\bibitem{inflation} J. di Giovanni, \c{S}. Kalemli-\"Ozcan, A. Silva and
M. A. Y\i ld\i r\i m. \emph{Pandemic-Era Inflation Drivers and Global
Spillovers}. Federal Reserve Bank of New York Staff Report 1080,
November 2023, revised July 2024. Consulted cover, abstract, introduction,
and sectoral supply-shock discussion, printed pp.~26--28, 9 October 2026.
\url{https://www.newyorkfed.org/medialibrary/media/research/staff_reports/sr1080.pdf}.
\end{thebibliography}

\end{document}
